% Pass 3 cuts (level of detail against Dablander's DPhil thesis) from thesis.tex header and Chapter 1 (Introduction).
% Compared with his Chapter 1 (thesis pp. 1-6): about three pages of motivation prose with no numbers, then Section 1.1 "Outline of Thesis and Research Contributions".
% That section has one paragraph per chapter: what it does, what it found (qualitative, no numbers), what is new, and where it is published.
% He has no research-questions section and no funding section in Chapter 1; funding (EPSRC, Lhasa) sits in his Acknowledgements.
% These notes planned more detail than his: result percentages and study counts in Section 1.4, lists of distances and representations, the Effiris citation and a blog-post count. Those are cut below. One line was added: each chapter paragraph states what is new.

% ===== FROM: Introduction notes, Paragraph 4 =====
% - Change: "vector operations" and "representations" become what was tried: Tanimoto on ECFP4, Euclidean on PDV, Mahalanobis; ECFP4, Sort \& Slice, PDV, canonical and randomised SMILES.

% ===== FROM: From noise robustness to adaptive representations, Paragraph 1 =====
% - Robustness is the normalised area under the R$^2$ retention curve as noise rises.
% - Under Gaussian noise on QM9, model architecture explains 49.5% of the variance in robustness; representation explains 9.2%.
% - (second half of a bullet) under clipping, models were most confident on the clipped labels.

% ===== FROM: From noise robustness to adaptive representations, Paragraph 2 =====
% - Twelve methods adapted to regression, 221 configurations, QM9 HOMO-LUMO gap with Gaussian noise.
% - All 20 configurations with the lowest R$^2$ retention judged a label by its neighbours' labels. This held for every distance tried.

% ===== FROM: Collaborations and funding, Paragraph 2 =====
% - Name their federated-learning platform, Effiris, in one clause \citep{Bassani2023, Hanser2023}.

% ===== FROM: Collaborations and funding, Paragraph 3 =====
% - (count) the collaboration produced five posts on its blog, Blopig.
